Suppose that @@M0@@. Then @@M1@@; thus @@M2@@ and @@M3@@. Moreover, @@M4@@; hence @@M5@@. Thus @@M6@@ is indeed a @@M7@@-sheaf. In addition, 6.5 (ii) shows that @@M8@@. If @@M9@@, we may suppose that @@M10@@ is a @@M11@@-sheaf and that @@M12@@, by induction on @@M13@@. Then 6.5 shows that @@M14@@ and @@M15@@; whence the assertion.

\sgaStableTarget{sga6.e13.thm.6-7}\textbf{Theorem 6.7.} \emph{There exist two sequences @@M16@@ and @@M17@@ of polynomials with the following properties. Assume that the conditions of 6.1 are satisfied and that @@M18@@ is a quotient of @@M19@@ with @@M20@@.}
\begin{enumerate}[label=(\roman*)]
\item \emph{If @@M21@@ is a @@M22@@-sheaf, then, for @@M23@@,}
@@M24@@
\item \emph{If @@M25@@ is the Hilbert polynomial of @@M26@@, then, for @@M27@@,}
@@M28@@
\end{enumerate}
\emph{(Notice that @@M29@@ and @@M30@@ since @@M31@@).}

Indeed, let an exact sequence @@M32@@ be such that @@M33@@ is a @@M34@@-sheaf; arguing by induction on @@M35@@, we may suppose that @@M36@@ have been defined, taking 1.7 into account. Then 6.4 (ii) makes the definition of @@M37@@ clear, taking 1.8 into account. Finally, (ii) follows formally from 6.4 (iii), 6.6, and assertion (i) (where we take @@M38@@ for @@M39@@).

